\documentclass[12pt,a4paper]{article}

\usepackage[T1]{fontenc}
\usepackage[utf8]{inputenc}
\usepackage[french]{babel}
\usepackage{geometry}
\usepackage{setspace}
\usepackage{hyperref}
\usepackage{booktabs}
\usepackage{longtable}
\usepackage{array}
\usepackage{enumitem}
\usepackage{graphicx}
\usepackage{xcolor}
\usepackage{fancyhdr}
\usepackage{titlesec}
\usepackage{amsmath}
\usepackage{amssymb}
\usepackage{listings}

\geometry{margin=2.5cm}
\onehalfspacing
\hypersetup{colorlinks=true, linkcolor=blue, urlcolor=blue, citecolor=blue}

\pagestyle{fancy}
\fancyhf{}
\lhead{SIMCO -- Rapport Technique}
\rhead{\thepage}

\titleformat{\section}{\large\bfseries}{\thesection}{0.6em}{}
\titleformat{\subsection}{\normalsize\bfseries}{\thesubsection}{0.6em}{}

\title{\textbf{SIMCO (Syst\`eme Intelligent Multimodal d'\'Evaluation Cognitive)\\Rapport Technique D\'etaill\'e}}
\author{Projet r\'ealis\'e par : Guy (MRGUY10)}
\date{24 mars 2026}

\begin{document}
\maketitle
\tableofcontents
\newpage

\section{R\'esum\'e du projet}
SIMCO est une plateforme d'\'evaluation cognitive multimodale. Le principe est simple :
\begin{enumerate}
    \item l'utilisateur passe un quiz,
    \item il d\'eclare son niveau de confiance,
    \item le syst\`eme observe en parall\`ele des signaux comportementaux,
    \item une estimation de \og confiance r\'eelle \fg{} est calcul\'ee,
    \item un rapport PDF est envoy\'e automatiquement par email.
\end{enumerate}

L'objectif principal est de ne pas seulement donner une note, mais aussi d'expliquer \emph{comment} la note s'aligne (ou non) avec la perception de l'apprenant.

\section{Pourquoi cette architecture ?}
Le projet est d\'ecoup\'e en plusieurs services pour trois raisons :
\begin{itemize}
    \item \textbf{Lisibilit\'e} : chaque service a une mission claire.
    \item \textbf{Maintenabilit\'e} : on corrige un service sans casser les autres.
    \item \textbf{Scalabilit\'e} : certains services lourds (vision, ML) peuvent \^etre d\'eploy\'es s\'epar\'ement.
\end{itemize}

\section{Architecture globale (diagramme simple)}
\subsection{Vue d'ensemble}
\begin{lstlisting}[basicstyle=\ttfamily\small]
                    +-----------------------+
                    |     Utilisateur       |
                    | (Navigateur / React)  |
                    +-----------+-----------+
                                |
                                | HTTP (quiz flow)
                                v
                    +-----------+-----------+
                    |   Quiz Backend        |
                    |   FastAPI :8000       |
                    +---+---------------+---+
                        |               |
      true-confidence   |               | notification email + PDF
      request           |               |
                        v               v
              +---------+-----+   +-----+----------------+
              | Confidence    |   | Notification Backend |
              | Backend :8010 |   | FastAPI :8020        |
              +---------------+   +----------+-----------+
                                                |
                                                | SMTP
                                                v
                                         +------+------+
                                         |   Serveur   |
                                         |    Mail     |
                                         +-------------+

 (Cote frontend) Webcam -> Face Backend :8084 pour la confiance visage
\end{lstlisting}

\subsection{R\^ole de chaque bloc}
\begin{longtable}{p{3.8cm} p{3.8cm} p{7.4cm}}
\toprule
\textbf{Service} & \textbf{Technologie} & \textbf{R\^ole fonctionnel} \\
\midrule
Frontend & React + Vite + Tailwind & Interface utilisateur, quiz, webcam, affichage des r\'esultats \\
Quiz Backend & FastAPI & Orchestrateur central (questions, r\'eponses, score, analyses, appels inter-services) \\
Face Backend & FastAPI + CV/ML & Analyse image/frame pour extraire des signaux visage/confiance \\
Confidence Backend & FastAPI + NumPy & Calcule une confiance \og r\'eelle \fg{} \`a partir des signaux agr\'eg\'es \\
Notification Backend & FastAPI + SMTP + ReportLab & G\'en\`ere le PDF de rapport et envoie l'email \\
\bottomrule
\end{longtable}

\section{Structure du d\'ep\^ot}
Les dossiers principaux sont :
\begin{itemize}
    \item \texttt{quiz-frontend/} : application React.
    \item \texttt{services/quiz\_backend/} : API principale.
    \item \texttt{services/face\_backend/} : API vision.
    \item \texttt{services/confidence\_backend/} : API logique de confiance.
    \item \texttt{services/notification\_backend/} : API email + PDF.
    \item \texttt{scripts/windows} et \texttt{scripts/unix} : scripts de d\'emarrage.
    \item \texttt{docs/} : documentation projet.
\end{itemize}

\section{Explication d\'etaill\'ee de chaque service}

\subsection{1) Frontend (\texttt{quiz-frontend})}
\subsubsection{Ce qu'il fait concr\`etement}
Le frontend est la porte d'entr\'ee utilisateur.
\begin{itemize}
    \item collecte le profil (nom, \^age, niveau, email obligatoire),
    \item lance un quiz personnalis\'e,
    \item envoie les r\'eponses,
    \item collecte la confiance globale de l'utilisateur,
    \item affiche les r\'esultats (score, histogrammes, analyses),
    \item affiche des alertes de focus en cours de quiz.
\end{itemize}

\subsubsection{Composants cl\'es}
\begin{itemize}
    \item \texttt{QuizPage.jsx} : logique principale du parcours quiz.
    \item \texttt{WebcamAnalyzer.jsx} : capture/traitement des m\'etriques webcam.
\end{itemize}

\subsubsection{Pourquoi c'est important}
Il traduit les donn\'ees techniques en informations compr\'ehensibles pour l'apprenant (feedback, interpr\'etation, plan d'action).

\subsection{2) Quiz Backend (\texttt{services/quiz\_backend})}
\subsubsection{Mission principale}
C'est le \textbf{chef d'orchestre} du syst\`eme.

\subsubsection{Ce qu'il fait pas \`a pas}
\begin{enumerate}
    \item Re\c{c}oit la demande de quiz et g\'en\`ere les questions via Ollama/Mistral.
    \item Cr\'ee une session utilisateur (id unique).
    \item Re\c{c}oit chaque r\'eponse et calcule le score en temps r\'eel.
    \item Stocke les donn\'ees comportementales envoy\'ees par le frontend.
    \item Re\c{c}oit la confiance globale d\'eclar\'ee en fin de quiz.
    \item Appelle le service de confiance (SIMCO Logic) pour estimer la confiance r\'eelle.
    \item Construit le r\'esultat final (question par question + recommandations).
    \item D\'eclenche l'envoi email via le service notification.
\end{enumerate}

\subsubsection{Endpoints principaux}
\begin{itemize}
    \item \texttt{POST /generate-quiz}
    \item \texttt{POST /submit-answer}
    \item \texttt{POST /update-confidence}
    \item \texttt{GET /quiz-results/\{session\_id\}}
\end{itemize}

\subsubsection{Logique de calibration (simplifi\'ee)}
La confiance d\'eclar\'ee est normalis\'ee sur $[0,1]$ puis compar\'ee au score observ\'e.

\[
C_n = \min(1,\max(0, C_{raw})) \text{ (si d\'ej\`a normalis\'e)}
\]
\[
C_n = \min(1,\max(0, C_{raw}/100)) \text{ (si exprim\'e en \%)}
\]

Un indice de d\'ecalage confiance/performance est utilis\'e pour classer le profil :
\[
DK\_index = \overline{C_{behavioral}} - Score_{\%}
\]

\subsubsection{Gestion d'erreurs}
\begin{itemize}
    \item Si \texttt{confidence\_backend} est indisponible : erreur contr\^ol\'ee c\^ot\'e API.
    \item Si \texttt{notification\_backend} \'echoue : le r\'esultat quiz est tout de m\^eme renvoy\'e (best effort).
\end{itemize}

\subsection{3) Face Backend (\texttt{services/face\_backend})}
\subsubsection{Mission principale}
Analyser les images webcam pour extraire des indices utilisables dans l'\'evaluation m\'etacognitive.

\subsubsection{Ce qu'il traite}
\begin{itemize}
    \item image unique (endpoint classify),
    \item lot de frames (endpoint classify/frames),
    \item calcul de pr\'edictions principales ou d\'etaill\'ees.
\end{itemize}

\subsubsection{Endpoints principaux}
\begin{itemize}
    \item \texttt{GET /api/health}
    \item \texttt{POST /api/classify}
    \item \texttt{POST /api/classify/frames}
\end{itemize}

\subsubsection{Point op\'erationnel critique}
Ce service d\'epend de biblioth\`eques lourdes (OpenCV, Keras, TensorFlow). Sous Windows, la compatibilit\'e est en pratique bien meilleure avec Python 3.10.

\subsection{4) Confidence Backend (\texttt{services/confidence\_backend})}
\subsubsection{Mission principale}
Transformer des signaux de confiance en une estimation plus robuste : la \textbf{confiance r\'eelle}.

\subsubsection{Entr\'ees}
\begin{itemize}
    \item \texttt{self\_confidence} (obligatoire, normalis\'ee dans $[0,1]$),
    \item \texttt{face\_confidence\_per\_question} (liste optionnelle de valeurs normalis\'ees).
\end{itemize}

\subsubsection{Sorties}
\begin{itemize}
    \item \texttt{true\_confidence\_normalized},
    \item \texttt{true\_confidence} en pourcentage,
    \item r\'esum\'e d'entr\'ee et type de mod\`ele.
\end{itemize}

\subsubsection{Endpoints principaux}
\begin{itemize}
    \item \texttt{GET /health}
    \item \texttt{POST /analyze/true-confidence}
\end{itemize}

\subsection{5) Notification Backend (\texttt{services/notification\_backend})}
\subsubsection{Mission principale}
Produire un rapport final lisible par un humain et l'envoyer automatiquement.

\subsubsection{Ce qu'il fait}
\begin{enumerate}
    \item Valide le payload (nom, email, r\'esultats, analyses).
    \item G\'en\`ere un PDF avec statistiques et graphiques.
    \item Envoie l'email via SMTP avec pi\`ece jointe.
\end{enumerate}

\subsubsection{Endpoints principaux}
\begin{itemize}
    \item \texttt{GET /health}
    \item \texttt{POST /notifications/quiz-result}
\end{itemize}

\subsubsection{Configuration}
Le service lit la configuration depuis \texttt{.env} \textbf{et} \texttt{.env.local}, ce qui \'evite des erreurs courantes de type \og credentials missing \fg{}.

\section{Flux complet d'un quiz (s\'equence pas \`a pas)}
\begin{lstlisting}[basicstyle=\ttfamily\small]
1. Frontend -> Quiz Backend      : POST /generate-quiz
2. Quiz Backend -> Ollama        : generation des questions
3. Frontend -> Quiz Backend      : POST /submit-answer (x N)
4. Frontend -> Quiz Backend      : POST /update-confidence
5. Quiz Backend -> Confidence API: POST /analyze/true-confidence
6. Quiz Backend -> Notification  : POST /notifications/quiz-result
7. Notification -> SMTP          : envoi email + PDF
8. Frontend <- Quiz Backend      : GET /quiz-results/{id}
\end{lstlisting}

\section{Ports et communication}
\begin{longtable}{p{6cm} p{6cm}}
\toprule
\textbf{Service} & \textbf{Port local} \\
\midrule
Quiz Backend & 8000 \\
Face Backend & 8084 \\
Confidence Backend & 8010 \\
Notification Backend & 8020 \\
Frontend (Vite) & 5173 \\
Ollama & 11434 \\
PostgreSQL & 5432 \\
\bottomrule
\end{longtable}

\section{Exploitation (run en local)}
\subsection{Scripts utiles}
\begin{itemize}
    \item \texttt{start-backend.bat}
    \item \texttt{start-face-class.bat}
    \item \texttt{start-simco-logic.bat}
    \item \texttt{start-notification.bat}
    \item \texttt{start-frontend.bat}
\end{itemize}

Ces wrappers appellent les scripts centralis\'es dans \texttt{scripts/windows/}.

\section{S\'ecurit\'e, limites et recommandations}
\subsection{Points \`a am\'eliorer en priorit\'e}
\begin{itemize}
    \item Ne jamais garder des secrets SMTP en clair en production.
    \item Ajouter authentification inter-services.
    \item Centraliser la journalisation et les m\'etriques.
    \item Ajouter retries/backoff sur les appels r\'eseau critiques.
\end{itemize}

\subsection{Risques techniques connus}
\begin{itemize}
    \item conflits de ports sous Windows,
    \item compatibilit\'e TensorFlow/Python pour le face backend,
    \item d\'ependances lourdes de la pile vision.
\end{itemize}

\section{Conclusion}
SIMCO est une architecture moderne et p\'edagogique :
\begin{itemize}
    \item le frontend rend l'exp\'erience claire,
    \item le quiz backend orchestre,
    \item le face backend observe,
    \item le confidence backend calibre,
    \item le notification backend restitue et diffuse.
\end{itemize}

Le syst\`eme est d\'ej\`a utilisable de bout en bout et la s\'eparation en services facilite fortement la compr\'ehension, la maintenance et l'\'evolution.

\appendix
\section{Commande de compilation LaTeX}
\begin{lstlisting}[basicstyle=\ttfamily\small]
pdflatex -output-directory=docs docs/SIMCO_Project_Report.tex
\end{lstlisting}

\end{document}
